\subsection{Preprocessing Boolean matrices for sparse operations}

Our algorithms exploit regularity to reduce dense BMM to a collection of somewhat sparse matrix multiplications. To this end, we need results on preprocessing matrices to speed up computations on sparse inputs. The first deals with multiplication of an arbitrary matrix with a sparse vector, and the second deals with multiplication of a sparse matrix with another (arbitrary) matrix.

\begin{theorem}[Blelloch-Vassilevska-Williams~\cite{ICALP08}] \label{processforsparse}
Let $B$ be a $n \times n$ Boolean matrix and let $w$ be the wordsize. Let $\kappa \geq 1$ and $\ell > \kappa$ be integer parameters. There is a data structure that can be constructed with $O((n^2 \kappa / \ell) \cdot \sum_{b=1}^{\kappa} \binom{\ell}{b})$ preprocessing time, so that for any Boolean vector $v$, the product $B \star v$ can be computed  in 
\[
O\left(n \log n + \frac{n^2}{\ell w} + \frac{n t}{\kappa w}\right)
\]
time, where $t$ is the number of nonzeros in $v$.
\end{theorem}

This result is typically applied as follows. Fix a value of $t$ to be the number of nonzeros we expect in a typical vector $v$. Choose $\ell$ and $\kappa$ such that $n/\ell \approx t/\kappa$, and $\sum_{b=1}^{\kappa} \binom{\ell}{b} = n^\delta$ for some $\delta > 0$. One such choice is $\kappa = \delta \ln(n)/\ln(en/t)$ and $\ell = \kappa \cdot en/t$ in which case we obtain:

\begin{theorem}
Let $B$ be a $n \times n$ Boolean matrix.  There is a data structure that can be constructed with $\tilde{O}(n^{2+\delta})$ preprocessing time, so that for any Boolean vector $v$, the product $B \star v$ can be computed in 
$$
O\left(n \log n + \frac{n t \ln(en/t)}{\delta w \ln n }\right)
$$
time, where $t$ is the number of nonzeros in $v$.
\end{theorem}

We should remark that we do not explicitly apply the above theorem, but the idea (of preprocessing for sparse vectors) is used liberally in this paper.

The following result is useful for multiplying a sparse matrix with another arbitrary matrix.
\begin{theorem}\label{sparsemm} There is an $O(m n \log(n^2/m)/(w\log n))$ time algorithm for computing $A \star B$, for every $n\times n$ $A$ and $B$, where $A$ has $m$ nonzeros and $B$ is arbitrary.
\end{theorem}
This result follows in a straightforward manner by combining the two lemmas below.
The first is a graph compression method due to Feder and Motwani.

\begin{lemma}[From Feder-Motwani~\cite{federmotwani}, Theorem 3.3]\label{sparseramsey} Let $\delta \in (0,1)$ be constant. We can write any $n \times n$ Boolean matrix $A$ with $m$ nonzeros as $A = (C \star D) \vee E$ where $C$, $D$ are $n \times m/n^{1-\delta}$, $m/n^{1-\delta} \times n$, respectively, both with at most $m (\log n^2/m)/(\delta \log n)$ nonzeros, and $E$ is $n \times n$ and has at most $n^{2-\delta}$ nonzeros. Furthermore, finding $C$, $D$, $E$ takes $O(mn^{\delta}\log^2 n)$ time.\end{lemma}

Since the lemma is not stated explicitly in~\cite{federmotwani}, let us sketch the proof for completeness. Using algorithmic Ramsey theoretic arguments (\ie, finding large bipartite cliques in a dense enough graph), Feder and Motwani show that for every bipartite graph $G$ on $2n$ nodes (with $n$ nodes each on left and right) and $m>n^{2-\delta}$ edges, its edge set can be decomposed into $m/n^{1-\delta}$ edge-disjoint bipartite cliques, where the total sum of vertices over all bipartite cliques (a vertex appearing in $K$ cliques is counted $K$ times) is at most $m (\log n^2/m)/(\delta \log n)$. Every $A$ can be written in the form $(C \star D) \vee E$, by having the columns of $C$ (and rows of $D$) correspond to the bipartite cliques. Set $C[i,k]=1$ iff the $i$th node of the LHS of $G$ is in the $k$th bipartite clique, and similarly set $D$ for the nodes on the RHS of $G$. Note that  $E$ is provided just in case $A$ turns out to be sparse.

We also need the following simple folklore result. It is stated in terms of wordsize $w$, but it can easily be implemented on other models such as pointer machines with $w=\log n$.

\begin{lemma}[Folklore]\label{sparseword} There is an $O(m n/w + pq + pn)$ time algorithm for computing $A \star B$, for every $p \times q$ matrix $A$ and $q \times n$  matrix $B$ where $A$ has $m$ nonzeros and $B$ is arbitrary.\end{lemma}

\begin{proof} We assume the nonzeros of $A$ are stored in a list structure; if not we construct this in $O(pq)$ time. Let $B_j$ be the $j$th row of $B$ and $C_i$ be the $i$th row of $C$ in the following. We start with an output matrix $C$ that is initially zero. For each nonzero entry $(i,j)$ of $A$, update $C_i$ to be the OR of $B_j$ and $C_i$. Each update takes only $O(n/w)$ time. It is easy to verify that the resulting $C$ is the matrix product.\end{proof}

\section{Combinatorial Boolean matrix multiplication via triangle removal}

In this section, we prove \expref{Theorem}{boolmmtriangles}.
That is, we show that a more efficient Triangle Removal Lemma implies more efficient Boolean matrix multiplication.
Let $A$ and $B$ be the matrices whose product $D$ we wish to compute. The key idea is to split the task into two cases. First, we use simple random sampling to determine the entries in the product that have many witnesses (where $k$ is a \emph{witness for $(i,j)$} if $A[i,k]=B[k,j]=1$). To compute the entries with few witnesses, we set up a tripartite graph corresponding to the remaining undetermined entries of the matrix product, and argue that it has few triangles. (Each triangle corresponds to a specific witness for a specific entry in $D$ that is still undetermined.) By a Triangle Removal Lemma, a sparse number of edges hit all the triangles in this graph.\footnote{Note that the triangle removal lemma may also return edges that do not lie in any triangle.} Using three carefully designed sparse matrix products (which only require one of the matrices to be sparse), we can recover all those entries $D[i,j]=1$ which have few witnesses.

We now describe our algorithm for BMM.

\smallskip

\noindent{\bf Algorithm:} Let $A$ and $B$ be $n \times n$ matrices. We wish to compute $D = A \star B$, \ie,
\[
D[i,j] = \left(\bigvee_{k=1}^n A[i,k] \wedge B[k,j]\right)\,.
\] 

\smallskip

\noindent
\emph{Random sampling for pairs with many witnesses.}
First, we detect the pairs $(i,j)$ with at least $\eps n$ witnesses.
Construct a $n\times n$ matrix $C$  as follows. Pick a sample $R$ of $(6\log n)/\eps$ elements from $[n]$.
For each $(i,j)$,  $1 \leq i,j \leq n$, check if there is a $k \in R$ that is a witness for $(i,j)$ in the product. If yes, set $C[i,j]=1$, otherwise $C[i,j]=0$. Clearly, this takes at most $O((n^2 \log n) /\eps)$ time. Note that  $C$ is dominated by the desired $D$, in that $C[i,j] \leq D[i,j]$ for all $i,j$.
If $(i,j)$ has at least $\eps n$ witnesses, then some witness lies in $R$ with probability at least $1-1/n^6$. Thus with probability at least $1-1/n^4$, $C[i,j]=D[i,j]=1$ for every $(i,j)$ with at least $\eps n$ witnesses.

\smallskip

\noindent
\emph{Triangle removal for pairs with few witnesses.} It suffices to determine those $(i,j)$ such that $C[i,j] = 0$ and $D[i,j] = 1$. We shall exploit the fact that such pairs do not have many witnesses. Make a tripartite graph $H$ with vertex sets $V_1$, $V_2$, $V_3$, each with $n$ nodes indexed by $1,\ldots,n$. Define edges as follows:
\begin{itemize}
\item Put an edge $(i,k) \in (V_1,V_2)$ if and only if $A[i,k]=1$.
\item Put an edge  $(k,j) \in (V_2,V_3)$  if and only if $B[k,j]=1$.
\item Put an edge $(i,j) \in (V_1,V_3)$ if and only if $C[i,j]=0$.\end{itemize}
That is, edges from $V_1$ to $V_3$ are given by $\overline{C}$, the complement of $C$. Observe that
$(i,k,j) \in (V_1,V_2,V_3)$ is a triangle if and only if $k$ is a witness for $(i,j)$ and $C[i,j] = 0$. Thus our goal is to find the pairs $(i,j)\in (V_1,V_3)$ that are in triangles of $H$.

Since every $(i,j) \in (V_1,V_3)$ has at most $\eps n$ witnesses, there are at most $\eps n^3$ triangles in $H$. Applying the promised Triangle Removal Lemma (in \expref{Theorem}{boolmmtriangles}), we can find in time $O(T(n))$ a set of edges $F$ where $|F| \leq f(\eps) n^2$ and each triangle must use an edge in $F$. Hence it suffices to compute those edges $(i,j) \in (V_1,V_3)$ that participate in a triangle with an edge in $F$. 

Define $A_F[i,j]=1$ if and only if $A[i,j]=1$ and $(i,j) \in F$. Similarly define $B_F$ and $\overline{C}_F$. Every triangle of $H$ passes through at least one edge from one of these three matrices. Let $T_A$ (resp. $T_B$ and $T_{\overline{C}}$) denote the set of triangles with an edge in ${A}_F$ (resp. $B_F$ and ${\overline{C}}_F$).  Note that we do not know these triangles.

We can determine the edges $(i,j) \in (V_1,V_3)$ that are in some triangle in $T_A$ or $T_B$ directly by computing $C_1 = {A}_F \star B$ and $C_2 = A \star {B}_F$, respectively. As $A_F$ and $B_F$ are sparse, by \expref{Theorem}{sparsemm}, these products can be computed in $O(|F|\log(n^2/|F|)/(w\log n))$ time. The $1$-entries of $\overline{C} \wedge C_1$ (resp. $\overline{C} \wedge C_2$) participate in a triangle in $T_A$ (resp. $T_B$). This determines the edges in $(V_1,V_3)$ participating in triangles from $T_A \cup T_B$.

Set $C = C \vee (C_1 \wedge \overline{C}) \vee (C_2 \wedge \overline{C})$, and update $\overline{C}$ and the edges in $(V_1,V_3)$ accordingly. The only remaining edges in $(V_1,V_3)$ that could be involved in a triangle are those corresponding to $1$-entries in ${\overline{C}}_F$. We now need to determine which of these actually lie in a triangle. 

Our remaining problem is the following: we have a tripartite graph on vertex set $(V_1,V_2,V_3)$ with at most $f(\eps) n^2$ edges between $V_1$ and $V_3$, and each such edge lies in at most  $\eps n$ triangles. We wish to determine the edges in $(V_1,V_3)$ that participate in triangles. This problem is solved by the following theorem.

\begin{theorem}[Reporting Edges in Triangles]\label{sparseedgelisting} Let $G$ be a tripartite graph on vertex set $(V_1,V_2,V_3)$ such that there are at most $\delta n^2$ edges in $(V_1,V_3)$, and every edge of $(V_1,V_3)$ is in at most $t$ triangles. Then the set of edges in $(V_1,V_3)$ that participate in triangles can be computed in $O(\delta n^3 \log(1/\delta)/(w \log n) + n^2 t)$ time. \end{theorem}

Setting $\delta = f(\eps)$ and $t = \eps n$, \expref{Theorem}{sparseedgelisting} implies the desired time bound in \expref{Theorem}{boolmmtriangles}. The idea of the proof of \expref{Theorem}{sparseedgelisting} is to work with a new tripartite graph where the vertices have asymptotically smaller degrees, at the cost of adding slightly more nodes. This is achieved by having some nodes in our new graph correspond to \emph{small subsets of nodes} in the original tripartite graph.

\begin{proof}[Proof of~\expref{Theorem}{sparseedgelisting}] We first describe how to do the computation on a pointer machine with $w=\log n$, then describe how to modify it to work for the word RAM.


\paragraph{Graph Construction} We start by defining a new tripartite graph $G'$ on vertex set $(V_1,V_2',V_3')$.
Let $\gamma < 1/2$. $V_2'$ is obtained by partitioning the nodes of $V_2$ into $n/(\gamma \log n)$ groups of size $\gamma \log n$ each. For each group, we replace it by $2^{\gamma \log n} = n^\gamma$ nodes, one corresponding to each subset of nodes in that group. Thus $V_2'$ has $n^{1+\gamma}/(\gamma \log n)$ nodes.

$V_3'$ is also constructed out of subsets of nodes. We form $n/\ell$ groups each consisting of $\ell$ nodes in $V_3$, where $\ell = \gamma (\log n) /(\delta \log(e/\delta))$. For each group, we replace it by $\binom{\ell}{k} \leq (el/k)^k \leq n^\gamma$ nodes, one corresponding to each subset of size up to $\kappa = \gamma (\log n)/(\log (e/\delta))$. So $V_3'$ has $O(n^{1+\gamma}/\ell)$ nodes.

\smallskip
\noindent
\emph{Edges in $(V_2',V_3')$:} Put an edge between $u$ in $V_2'$ and $x$ in $V_3'$ if there is an edge $(i,j)$ in $(V_2,V_3)$ such that $i$ lies in the set corresponding to $u$, and $j$ lies in the set corresponding to $x$. For each such edge $(u,x)$, we make a list of all edges $(i,j) \in (V_2,V_3)$ corresponding to it. Observe the list for a single edge has size at most $ \gamma \log n \cdot \ell = O(\log^2 n)$. 

\smallskip
\noindent
\emph{Edges in $(V_1,V_2')$:} The edges from $v \in V_1$ to $V_2'$ are defined as follows. For each group in $V_2$ consider the neighbors of $v$ in that group. Put an edge from $v$ to the node in $V'_2$ corresponding to this subset. Each $v$ has at most $n/(\gamma \log n)$ edges to nodes in $V_2'$.

\smallskip
\noindent
\emph{Edges in $(V_1,V_3')$:} Let $v \in V_1$. For each group $g$ of $\ell$ nodes in $V_3$, let $N_{v,g}$ be the set of neighbors of $v$ in $g$. Let $d_{v,g} = |N_{v,g}|$. Partition $N_{v,g}$ arbitrarily into $t=\lceil d_{v,g}/\kappa \rceil$ subsets $s_1,\ldots,s_t$ each of size at most $\kappa$. Put edges from $v$ to $s_1,\ldots,s_t$ in $V_3'$. The number of these edges from $v$ is at most $\sum_g \lceil d_{v,g}/\kappa \rceil \leq n/\ell + d_v /\kappa$, where $d_v$ is the number of edges from $v$ to $V_3$. Since $\sum_v d_v \leq \delta n^2$, the total number of edges from $V_1$ to $V_3'$ is $O(\delta \log (1/\delta) n^2 /(\gamma \log n))$.

\paragraph{Final Algorithm} For each vertex $v \in V_1$, iterate over each pair of $v$'s neighbors $u \in V_2'$ and $x \in V_3'$. If $(u,x)$ is an edge in $G'$, output the list of edges  $(i,j)$ in $(V_2,V_3)$ corresponding to $(u,x)$, otherwise continue to the next pair. From these outputs we can easily determine the edges $(v,j)$ in $(V_1,V_3)$
that are in triangles: $(v,j)$ is in a triangle if and only if node $j$ in $V_3$ is output as an end point of some edge $(i,j) \in (V_2,V_3)$ during the loop for $v$ in $V_1$.

\smallskip
\noindent
\emph{Running Time:} The graph construction takes at most $O(n^{2+2\gamma})$. In the final algorithm, the total number of pairs $(u,w)$ in $(V_2',V_3')$ that are examined is at most \[(n/\log n) \cdot O(\delta n^2 (\log 1/\delta)/ \log n)) \leq O(\delta \log (1/\delta) n^3/\log^2 n)\,.\]

We claim that the time used to output the lists of edges is at most $O(n^2 t)$ time. A node $j$ from $V_3$ is on an output list during the loop for $v$ in $V_1$ if and only if $(v,j)$ is an edge in a triangle, with some node in $V_2$ that has a 1 in the node $i$ in $V_2'$. Since each edge from $(V_1,V_3)$ in a triangle is guaranteed to have at most $t$ witnesses in $V_2$, the node $j$ is output at most $t$ times over the loop for $v$ in $V_1$. Hence the length of all lists output during the loop for $v$ is at most $nt$, and the total time for output is at most $O(n^2 t)$.

\smallskip
\noindent
\emph{Modification for $w$-word RAM:} We now show how to replace a $\log$-speedup by a $w$-speedup with wordsize $w$.
We form $V_3'$ as above and consider the graph on $(V_1,V_2,V_3')$ (note the $V_2$ instead of $V_2'$ previously).
Recall that a node $x \in V_3'$ corresponds to a subset $V_x \subset V_3$ of at most $\kappa$ vertices (from some group of size $\ell$).
An edge $(v,x) \in (V_1,V_3')$ implies that $v$ is adjacent to every vertex in $V_x$.
For each $v \in V_1$, let $S_v$ be the $n$-bit indicator vector corresponding to neighbors of $v$ in $V_2$. For each
$x \in V_3'$, let $T_x$ denote the $n$-bit indicator vector corresponding to the union of neighbors of vertices $V_x$ in $V_2$, \ie, $T_x[i]=1$  iff some $v \in V_x$ is adjacent to $i\in V_2$.
The vectors $S_v$ and $T_x$ are stored as $n/w$ words.  With each bit $T_x[i]$ of $T_x$ such that $T_x[i]=1$, we store a pointer to a list of vertices $v \in V_x$  that are adjacent to  $i \in V_2$. It is easily checked that the overall space usage is bounded by  $O(n \cdot |V_3'| \cdot \kappa) = \tilde{O}(n^{2+\gamma})$.

The algorithms works as follows: Now for every $v$ in $V_1$ and every neighbor $x \in V'_3$ of $v$, for $i=1,\ldots,n/w$, we look up
the $i$th word $q$ from $S_v$ and the $i$th word $q'$ in the $T_x$, and compute $q \wedge q'$. If this is nonzero, then each bit location $b$ where $q\wedge q'$ has a 1 means that the node corresponding to $b$ forms a triangle with $v$ and some vertex in $V_x$. For each such bit (say at location $i$), the list associated to it consists of precisely the nodes in $V_x$ forming a triangle with $i \in V_2$ and $v$.

The running time follows as there are $O(\delta n^2 (\log 1/\delta)/ \log n))$ edges in $(V_1,V_3')$, and hence there we look up at most $O(n (\delta n^2 (\log 1/\delta)/ \log n)  (n/w)$ words $q,q'$. Since the total number of triangles is $n^2 t$, the total time spent in processing whenever
 $q \wedge q'=1$ is $O(n^2 t)$.
\end{proof}


\begin{remark} Note that we only use randomness in the BMM algorithm  to determine the pairs $(i,j)$ that have many  witnesses.  Moreover, by choosing a larger sample $R$ in the random sampling step (notice we have a lot of slack in the running time of the random sampling step), the probability of failure can be made exponentially small.\end{remark}
